\begin{figure}[htbp]
\centering
\begin{tikzpicture}
\begin{axis}[
    ybar,
    width=0.9\textwidth,
    height=7cm,
    ylabel={Покрытие уровней Блума, \%},
    ymin=0, ymax=100,
    xtick=data,
    symbolic x coords={Нормативные,{Все документы},{Учебные программы},Продуктовые},
    x tick label style={font=\small, rotate=25, anchor=east},
    bar width=22pt,
    enlarge x limits=0.2,
    nodes near coords,
    nodes near coords style={font=\small, anchor=south},
]
\addplot[fill=red!50, draw=red!70!black] coordinates
    {(Нормативные,40) ({Все документы},61.5) ({Учебные программы},66.7) (Продуктовые,75)};

% Целевой порог
\draw[dashed, thick, black!70] (axis cs:Нормативные,70) -- (axis cs:Продуктовые,70)
    node[pos=0.75, above, font=\scriptsize] {Целевой порог (70\,\%)};

% Перекрасить Продуктовые (>=70) — наложим поверх
\addplot[fill=teal!70!black, draw=teal!50!black, forget plot] coordinates
    {(Нормативные,0) ({Все документы},0) ({Учебные программы},0) (Продуктовые,75)};
\end{axis}
\end{tikzpicture}
\caption{Покрытие уровней таксономии Блума по категориям документов (конфигурация B2)}
\label{fig:bloom-by-category}
\end{figure}
